This study closes with the version dated 13 October 2026. The following conclusions do not establish what happened or why. They describe what the results of a numerical model suggest, given the input parameters declared in this document and the 84 completed simulations: 19 screening runs with the previous terrain, 41 with the full city, 5 concentration quadrature runs and 19 campus micro-simulations. With different input data, another model or new observations, the results may change.

\subsection*{What the Results Suggest}
\begin{enumerate}
\item \textbf{General structure.} With the full city and the assumptions adopted, the model reproduces the Camino El Alba, General Blanche and Quebrada Honda corridors that delimit the official affected area, and the reach to Padre Hurtado. It does not reproduce the spatial distribution of impacts within that area: the test using held-out observations, defined before examining the results and performed only once, had an unfavorable outcome (19 of 49 facilities, compared with the specified \SI{70}{\percent}).
\item \textbf{Volume.} In the scenarios tested, about \SI{300000}{m^3} of mixture allows the observed reach to be reproduced. This is the same order of magnitude as an independent rainfall-balance estimate adopting the reported soil saturation. Neither value is a measurement.
\item \textbf{Flow inputs to the model.} Scenarios in which mud emerges at a single point do not reproduce the full set of observations. The most compatible scenarios combine inflow through the main channel, which reproduces the northern branch in the model, with the southern branch and assumed hillside rainfall, which reproduce part of the reported impacts in the southern sector.
\item \textbf{Role of the city.} Representing the 42,018 buildings changes the result: in the model, streets concentrate the flow and increase depths.
\item \textbf{Drainage.} The capacity scenarios for the Francisco Bulnes Correa culvert are hypothetical. In the model, with the assumed volume, a calculated discharge of the order of \SI{65}{m^3/s} reaches the inlet. With lower assumed capacities, flow continues over the surface along Quebrada Honda. These scenarios do not describe the real structure, for which public information on capacity, cross section and condition was unavailable. They do not support claims that the structure was obstructed, that it malfunctioned or that a different capacity would have changed what happened.
\item \textbf{Universidad de los Andes campus.} In the simulations, the assumed turbid-water inflow from the hillside and hill follows La Plaza Sur and enters a sunken courtyard through a clearing; this calculated route is compatible with the images consulted. With a fine mesh (m36), the calculated depth in the courtyard reaches \SI{0.17}{m}, within the range estimated from the published video, and La Plaza Sur is wetted at \SI{58}{\percent} of its points. The simulated overflow from Álvaro del Portillo does not reach that courtyard with this terrain.
\item \textbf{Fire.} According to the severity estimated with Sentinel-2, \SI{82}{\percent} of the contributing catchments shows a burn signal and \SI{10}{\percent} has moderate-high or high severity. With the assumed coefficients, the balance without fire yields about \SI{62}{\percent} of the mud volume and \SI{43}{\percent} of the solids of the scenario with fire. Staley M1, used illustratively, yields a median probability of 0.28 with fire versus 0.07 without it. In the city, the corrected runs without fire (f35b and f38b) wet \SIrange{179}{193}{ha}, compared with \SI{222}{ha} in p34b (\SIrange{13}{19}{\percent} less), and calculate \SIrange{0.19}{0.23}{m} next to Clínica UC Christus. The ECI remains at 80--82, compared with 81 with fire, and does not distinguish these scenarios.
\item \textbf{Best fit.} The integrated hypothesis obtains an exploratory compatibility index (ECI) of 79--82/100: 82 with an assumed \SI{35}{mm/h} and a fine time step (z32), 81 with the pulse measured at Río Molina (p34b, corrected) and 79 with that rainfall distributed over \SI{4}{h} (v33). It combines distributed channels, rainfall on the lower hillside, houses that allow mud to enter, turbid water on the lower hillside and campus features. It is a fit made after the held-out test: it is not validated and the index is not a measure of accuracy.
\item \textbf{Hypotheses unsupported by the model.} In the tests, carving the \emph{bike park} trail does not change the calculated water route, and an assumed log obstacle at the end of the tree line does not divert the simulated flow southward.
\end{enumerate}

\subsection*{What Remains Unproven}
\begin{itemize}
\item \textbf{Independent validation.} No hypothesis developed after the held-out test was compared with new observations, such as Sentinel-1 radar, high-water marks or depths measured in the field.
\item \textbf{Rainfall in the catchment.} No rain gauge within the catchment was identified in the sources consulted. The Río Molina station, about \SI{7}{km} away, was used with an assumed time offset.
\item \textbf{Actual flow volume and composition.} No measurements of the mobilized volume, its concentration or its rheology were located in the sources consulted. Alternatives to the Beta(3,2) distribution used in the quadrature, and the estimation of inundation probabilities and extreme percentiles, remain untested.
\item \textbf{Drainage structures.} The actual capacity, cross section and condition of the culvert and conduits are unknown. The study's scenarios do not evaluate them.
\item \textbf{Clínica UC Christus.} No configuration reproduces both the reported depth of about \SI{1}{m} \citep{cadena3_2026} and the extent within the official affected area. Calibration of retention in homes using damage information by dwelling remains untested.
\item \textbf{Terrain.} The \SI{30}{m} elevation model does not resolve roadways, curbs or narrow passages. The carved features are coarse representations based on supplied observations. No LiDAR terrain model was available.
\item \textbf{Processes not simulated:} erosion and deposition along the flow path, entrained logs and vehicles, urban drainage, infiltration within the city and changes in flow concentration. The tracer that would separate hydrograph inflow from rainfall, the sediment hydrograph and concentration transport remain untested.
\item \textbf{Fire parameters.} Runoff increases by severity are assumed; their local calibration and the time of concentration of each subcatchment by severity remain untested. Severity thresholds and empirical volume and probability models were calibrated outside Chile.
\item \textbf{Hypothetical tree obstacle.} Only one assumed location and height were tested, on the available terrain. It was not verified in the field.
\end{itemize}

\subsection*{Nature and Scope of This Study}
This document is an independent technical exercise for academic and public-information purposes, undertaken on personal initiative and not commissioned by third parties. It is not an expert assessment, an expert report, an official technical opinion or an evaluation of the structures, services or decisions of any person or institution. It is based exclusively on public information consulted between 8 and 13 October 2026, recorded with dates and hashes in the evidence log. The author has no direct knowledge of the events beyond these sources.

The results are responses of a numerical model to declared input parameters and explicit assumptions, several of them hypothetical. No calculated figure is a measurement, and no statement in this document asserts the cause of the event or attributes responsibility. References to structures, services or institutions concern model scenarios or what the cited sources reported, with attribution. The sources disagree on several points, and those discrepancies are presented without resolving them.

The study is closed as of this version. The code, derived data and frozen versions are available for review and reproduction. Any use of these results in technical, administrative or judicial decisions requires an independent professional assessment, with design information, inspection and field measurements.
